\section{Why countably many Radon probes cannot recover full original capacity}
\label{sec:r12-countable-measure-capacity-obstruction}

The exact energy-normalized capacity in the preceding section is
controlled by one original Radon measure. Its null ideal is therefore
different in nature from the full original tester capacity. The
following all-dimensional example proves that the distinction cannot
in general be removed by replacing that measure by a countable family.
This is a counterexample within the allowed background class, not a
restriction on the backgrounds of the positive theorems.

Fix integers $n\geq q\geq1$ and a real number $L>0$. Work over the
original finite extension $K/\mathbf Q_p$ with its unchanged absolute
value. On $X=(\mathbf G_{m,K}^n)^{\rm an}$ put
\[
 \begin{gathered}
 t_i=\log|T_i|,\qquad
 V=\{x:-L<t_i(x)<L\text{ for }1\leq i\leq n\},\\
 Y=\{T_{q+1}=\cdots=T_n=1\},\qquad R=[Y]|_V.
 \end{gathered}
 \tag{CM1}
\]
If $q=n$, $Y=X$. The current $R$ is an original symmetric strongly
positive, hence middle-positive, separately closed current of degree
$(n-q,n-q)$. It extends across the closed box. Taking a larger compact
polyannulus supplies the precise strict-box hypothesis. In particular
no original field, absolute value or boundary convention has changed.

For $a\in(-L,L)^q$, denote by $x_a$ the original Gauss point of $Y$
with radii $(e^{a_1},\ldots,e^{a_q})$. These points exist for arbitrary
real $a$, without a value-group or algebraic-closure assumption.
The map $a\mapsto x_a$ is a continuous injection. Define on the
whole original domain $V$
\[
                      h_a=\frac{1}{2L}
                                  \max_{1\leq i\leq q}|t_i-a_i|.
 \tag{CM2}
\]

\begin{lemma}[A full original normalized tester at every Gauss point]
\label{lem:r12-gauss-singleton-full-capacity}
The function $h_a$ is an original continuous classical psh function
with $0\leq h_a\leq1$ on $V$. Its full original top current is
\[
                         (Dh_a)^q\wedge R=L^{-q}\delta_{x_a}.
 \tag{CM3}
\]
Consequently
\[
             C_R(\{x_a\};V)\geq L^{-q},\qquad
             C_R^*(\{x_a\};V)\geq L^{-q}.
 \tag{CM4}
\]
These are assertions for the unrestricted original normalized tester
family on the whole $V$, not for a fixed finite presentation class.
\end{lemma}

\noindent\textbf{Status: \texttt{DIRECT\_COROLLARY}} of the
frozen original bounded BT calculus, its integration-current
projection formula, and the elementary toric convex PL calibration.
The explicit calibration below fixes its normalization. It is not a
new instance of the arbitrary-background analytic relation quotient.
The exact frozen torus locator is R6
\path{sections/capacity_gauss_convex_reduction.tex},
\texttt{eq:r6-toric-convex-laplacian}, whose source SHA256 is
\texttt{2ab2277e81fdad829815f0fe30dc6e921f1d0d8c1fb70de71b718fe4f9d83e38}.
It explicitly distinguishes the all-original-test toric identity
from a mere pushforward identity. The original real-radius
calibration is also recorded in R6
\path{sections/capacity_laminar_graphs.tex},
\texttt{eq:r6-arbitrary-radius-capacity}. The targeted nested
archive and source checks are retained in
\path{checks/targeted_r6_calibration.json}.

\begin{proof}
The function in (CM2) is the pullback of a convex real PL function
under an actual original unit moment. Convex smoothing and the frozen
bounded uniform convergence theorem make it original classical psh.
For all $x\in V$, $|t_i(x)-a_i|<2L$, which proves its normalization.

On the coordinate subtorus $Y$, the subgradient at $a$ is the
cross-polytope
\[
                  P_a=\operatorname{conv}
                   \left\{\frac{e_i}{2L},-\frac{e_i}{2L}:1\leq i\leq q\right\},
             \qquad \operatorname{vol}(P_a)=\frac{1}{q!L^q}.
 \tag{CM5}
\]
Here is a direct smooth calculation that also checks support on
the full original test sheaf. Put
\[
 F(z)=\log\!\left(\sum_{i=1}^q
              \bigl(e^{z_i/(2L)}+e^{-z_i/(2L)}\bigr)\right),
       \qquad g_\epsilon(t)=\epsilon F((t-a)/\epsilon).
 \tag{CM5a}
\]
Then $0\leq g_\epsilon-h_a\leq\epsilon\log(2q)$ on the entire
real moment cube. The Hessian of $F$ is the positive definite
covariance matrix of the affinely spanning vectors
$\{\pm e_i/(2L)\}$ with strictly positive exponential weights.
The map $\nabla F$ is a diffeomorphism onto the interior of $P_a$:
for any interior $v$, the strictly convex function $F(z)-v\cdot z$
is coercive, since the support function of $P_a$ exceeds $v\cdot z$
by a positive multiple of $|z|$. It has one critical point;
the positive Hessian gives the smooth inverse. Change of variables
therefore gives
\[
                   \int_{\mathbf R^q}\det(\nabla^2F(z))\,dz
                          =\operatorname{vol}(P_a).
 \tag{CM5b}
\]
Since $\det\nabla^2g_\epsilon(t)
=\epsilon^{-q}\det\nabla^2F((t-a)/\epsilon)$, its finite
measure converges weakly to $\operatorname{vol}(P_a)\delta_a$.
This follows directly by dominated convergence after the displayed
change of scale; no general Monge--Ampere continuity theorem for
unrestricted unbounded potentials is involved.

For completeness the elementary original torus integration formula
has its full-test meaning here. At any point $x$ of the $q$-torus
outside its Gauss skeleton, some Laurent polynomial $f=\sum_m c_mT^m$
satisfies
$|f(x)|<\max_m |c_m|e^{m\cdot t(x)}$. This strict inequality persists
on an ordinary neighborhood. At each point of that neighborhood at
least two monomials attain the displayed maximum, by the ultrametric
inequality. Its moment image thus lies in the tropical hypersurface
of $f$, of real dimension at most $q-1$. Every smooth moment-pulled
$(q,q)$ form vanishes there by the definition of original CLD pullback.
The resulting smooth top measure is therefore carried by the original
Gauss skeleton. On that skeleton canonical torus tropical integration
has weight one and gives, for every compact original smooth scalar
test $\eta$,
\[
  \int_Y \eta\,(Dg_\epsilon)^q
       =q!\int_{(-L,L)^q}\eta(x_t)
                          \det\nabla^2g_\epsilon(t)\,dt.
 \tag{CM5c}
\]
Skeleton support is what makes this an identity on all original
tests, not merely a statement about the moment pushforward.
The normalization agrees with the unit simplex radial shell (PO6).
Now (CM5b), continuity of $t\mapsto\eta(x_t)$ and the uniform
bounded-potential convergence theorem give
$q!\operatorname{vol}(P_a)\eta(x_a)=L^{-q}\eta(x_a)$.
For $q<n$, the frozen integration-current projection formula applies
this same calculation on $Y$ and pushes the Dirac measure into the
original ambient $V$. Thus all original tests and all actual
moment refinements are retained in (CM3).
Using $h_a$ as a permitted normalized tester gives the first
inequality in (CM4); monotonicity of the ordinary-open outer
regularization gives the second.
\end{proof}

\begin{theorem}[No countable Radon family controls the full capacity-null ideal]
\label{thm:r12-no-countable-radon-capacity-control}
For the allowed data (CM1), let $(\nu_j)_{j\geq1}$ be any countable
family of positive Radon measures on $V$. There is one original
point $x_a$, lying in a fixed ordinary relatively compact window
independent of the measures, such that
\[
            \nu_j(\{x_a\})=0\quad\hbox{for every }j,
                  \qquad C_R^*(\{x_a\};V)\geq L^{-q}>0.
 \tag{CM6}
\]
Hence the full original outer-capacity-null ideal cannot be
characterized as the simultaneous null ideal of countably many
positive Radon measures. In particular, there is no single Radon
measure $\nu$ and function $F:[0,\infty)\to[0,\infty)$ with
$F(0)=0$ such that
\[
                         C_R^*(A;V)\leq F(\nu(A))
 \tag{CM7}
\]
for every Borel subset of that window.
\end{theorem}

\noindent\textbf{Status: \texttt{NEW\_PROVED}} as an
all-dimensional obstruction to the proposed countable Radon
control mechanism. The proposed universal equivalence and (CM7)
are \texttt{FALSE}.

\begin{proof}
Choose the fixed window
$W=\{-3L/4<t_i<3L/4\ (1\leq i\leq n)\}\Subset V$
and the compact uncountable Gauss set
\[
                    K_G=\{x_a:a\in[-L/2,L/2]^q\}\subset W.
 \tag{CM8}
\]
Each $\nu_j$ is finite on $K_G$. A finite measure has at most
countably many atoms there: for each positive integer $m$, only
finitely many atoms have mass at least $1/m$. Thus the union,
over all $j$, of the atoms of $\nu_j$ in $K_G$ is countable.
Choose $x_a$ outside it. All the vanishing assertions in (CM6)
follow, while (CM4) gives the strictly positive full capacity.
Applying the proposed (CM7) to this singleton would give
$L^{-q}\leq F(0)=0$, a contradiction.
\end{proof}

\begin{corollary}[Strict null-ideal separation for the parameter energy capacity]
\label{cor:r12-parameter-capacity-strict-null-ideal}
In every $n\geq q\geq1$, choose (CM1) and use the empty parameter
tuple in the parameter finite-energy construction. For any one of
its allowed probes $\psi$, there is a Borel singleton $A\subset W$
such that
\[
                  \mathfrak c_\psi(A)
                    =\mathfrak c_\psi^{\rm o}(A)=0,
                         \qquad C_R^*(A;V)>0.
 \tag{CM9}
\]
The same singleton can be chosen simultaneously for any countable
family of such probes and their energy-normalized capacities.
\end{corollary}

\noindent\textbf{Status: \texttt{NEW\_PROVED}} explicit
strict separation, using the already proved exact null ideal
in (PC12).

\begin{proof}
The empty parameter tuple makes conormality automatic; it does
not change the current $R$. Apply (CM6) to the finite positive
Radon measure $\omega=(D\psi)^q\wedge R$, or to the countable
family of these measures. Formula (PC12) gives precisely the
two zeros in (CM9), whereas (CM4) gives its last inequality.
All resulting capacities still use their explicitly different
admissible families. No comparison of those families has been
suppressed.
\end{proof}

\paragraph{What this does not contradict.}
The geometry-only probe detects specified normal or Jensen defects
by closedness and two-sided original test perturbations; it was
never asserted to dominate every normalized source measure.
Countable cofinality of original smooth test stages is likewise
not countable cofinality of all normalized nonlinear tester
\emph{measures}. The band formulas choose their essential
suprema for each fixed set and parameter control; they do not
claim a single countable list detecting every analytic set.
Thus (CM6) is an obstruction to a particular route to unrestricted
capacity, not a counterexample to its Choquet regularity or to
the analytic PS/PL relation quotient.
